\section{用例}
\subsection{フォント}

\subsubsection{書体}

\ctable[pos=H,left]{llll}{}{
  \FL
  & 書体
  & 欧文
  & 和文
  \ML
  & Roman
  & \textrm{0123456789ABCDEFGHIJKLMNOPQRabcdefghijklmn}
  & \textrm{あいうえお}
  \NN
  & San Serif
  & \textsf{0123456789ABCDEFGHIJKLMNOPQRabcdefghijklmn}
  & \textsf{あいうえお}
  \NN
  & Type Writer
  & \texttt{0123456789ABCDEFGHIJKLMNOPQRabcdefghijklmn}
  & \texttt{あいうえお}
  \NN
  & 明朝体
  & \textmc{0123456789ABCDEFGHIJKLMNOPQRabcdefghijklmn}
  & \textmc{あいうえお}
  \NN
  & ゴシック体
  & \textgt{0123456789ABCDEFGHIJKLMNOPQRabcdefghijklmn}
  & \textgt{あいうえお}
  \ML
  & ミディアム
  & \textmd{0123456789ABCDEFGHIJKLMNOPQRabcdefghijklmn}
  & \textmd{あいうえお}
  \NN
  & ボールド
  & \textbf{0123456789ABCDEFGHIJKLMNOPQRabcdefghijklmn}
  & \textbf{あいうえお}
  \NN
  & イタリック
  & \textit{0123456789ABCDEFGHIJKLMNOPQRabcdefghijklmn}
  & \textit{あいうえお}
  \LL
}

\subsubsection{deluxeフォント}

\ctable[pos=H,left]{llll}{}{
  \FL
  & 書体 & フォント & テキスト
  \ML
  & 明朝体 & ヒラギノ明朝Ｗ３
  & \mcfamily{永遠の美を追求したタイポグラフィー}
  \NN
  & 明朝太字 & ヒラギノ明朝Ｗ６
  & \mcfamily\bfseries{永遠の美を追求したタイポグラフィー}
  \NN
  & ゴシック & ヒラギノ角ゴＷ３
  & \gtfamily{永遠の美を追求したタイポグラフィー}
  \NN
  & ゴシック太字 & ヒラギノ角ゴＷ６
  & \gtfamily\bfseries{永遠の美を追求したタイポグラフィー}
  \NN
  & ゴシック極太 & ヒラギノ角ゴＷ８
  & \gtfamily\ebseries{永遠の美を追求したタイポグラフィー}
  \NN
  & 丸文字 & ヒラギノ丸ゴＷ４
  & \mgfamily{永遠の美を追求したタイポグラフィー}
  \NN
  & 明朝細字 & ヒラギノ明朝Ｗ２
  & \mcfamily\ltseries{永遠の美を追求したタイポグラフィー}
  \LL
}

\clearpage

\subsubsection{文字修飾}
\begin{itemize}
\item 普通の文字
\item \emph{強調(emph)}
\item \sout{打ち消し線}
\item $\overline{上線}$
\item $^{上付き}$
\item $_{下付き}$
\end{itemize}

\subsubsection{フォントサイズ}
\def\sampletext{Windowsでコンピュータの世界が広がります。}
\begin{itemize}
  \normalsize\item\tiny{\sampletext(tiny)}
  \normalsize\item\scriptsize{\sampletext(scriptsize)}
  \normalsize\item\footnotesize{\sampletext(footnotesize)}
  \normalsize\item\small{\sampletext(small)}
  \normalsize\item\normalsize{\sampletext(normalsize)}
  \normalsize\item\large{\sampletext(large)}
  \normalsize\item\Large{\sampletext(Large)}
  \normalsize\item\LARGE{\sampletext(LARGE)}
  \normalsize\item\huge{\sampletext(huge)}
  \normalsize\item\Huge{\sampletext(Huge)}
  \normalsize\item\HUGE{\sampletext(HUGE)}
\end{itemize}
\normalsize

\subsection{リンク}
\url{http://localhost/}

\href{http://www.example.com/}{example.com}

\subsection{quote環境}

\begin{quote}
吾輩は猫である。名前はまだ無い。

どこで生れたかとんと見当がつかぬ。
\end{quote}

\subsection{verbatim環境}

半角80文字がちょうど収まるサイズにしています。
コマンドラインやソースコードを示す場合はverbatim環境よりもlistingsを推奨します

\begin{verbatim}
01234567890123456789012345678901234567890123456789012345678901234567890123456789
０１２３４５６７８９０１２３４５６７８９０１２３４５６７８９０１２３４５６７８９
\end{verbatim}

フォントが小さくなって嫌だ、という方のために60文字がちょうど収まるverbatim60とverbatim40を用意しています。

\subsection{verbatim60}

\begin{verbatim60}
012345678901234567890123456789012345678901234567890123456789
０１２３４５６７８９０１２３４５６７８９０１２３４５６７８９
\end{verbatim60}

\subsection{verbatim40}

\begin{verbatim40}
0123456789012345678901234567890123456789
０１２３４５６７８９０１２３４５６７８９
\end{verbatim40}

\subsection{コードリスティング}

\subsubsection{コマンドライン}
\begin{lstlisting}[language=command]
# uname --machine
x86_64
\end{lstlisting}

\subsubsection{C}
\begin{lstlisting}[language=C]
#include <stdio.h>

int main(int argc, char *argv[]){
    printf("Hello World!");
    return 0;
}
\end{lstlisting}

\subsubsection{Python}
\begin{lstlisting}[language=python]
def hello():
    print('Hello World!')

if __name__ == '__main__':
    hello()
\end{lstlisting}

\subsubsection{80文字テスト}
\begin{lstlisting}
01234567890123456789012345678901234567890123456789012345678901234567890123456789
\end{lstlisting}

\subsection{リスト}

\subsubsection{列挙リスト}
\begin{itemize}
\item りんご
\item みかん
\item バナナ
\end{itemize}

\subsubsection{順序リスト}
\begin{enumerate}
\item 睦月
\item 如月
\item 弥生
\end{enumerate}

\subsubsection{定義リスト}
\begin{description}
\item[Monday] 月曜日
\item[Tuesday] 火曜日
\item[Wednesday] 水曜日
\end{description}

\subsection{テーブル}

\subsubsection{普通のテーブル}
\begin{table}[htbp]
\begin{center}
\begin{tabular}{|l|c|r|}
\hline
セル1 & セル2 & セル3 \\ \hline
セル4 & セル5 & セル6 \\ \hline
セル7 & セル8 & セル9 \\ \hline
\end{tabular}
\caption{キャプション}
\end{center}
\end{table}

\subsubsection{longtable}
\begin{longtable}[c]{@{}crr@{}}
\toprule\addlinespace
果物 & 値段 & 個数
\\\addlinespace
\midrule\endhead
バナナ & 134円 & 10本
\\\addlinespace
オレンジ & 210円 & 10個
\\\addlinespace
\bottomrule
\end{longtable}

\subsection{水平線}
\begin{center}\rule{\textwidth}{.5pt}\end{center}

\subsection{画像}
\includegraphics[width=0.8\textwidth]{latex.eps}

\clearpage
\subsection{数式}
\begin{enumerate}

\item 総和
\begin{equation}
\sum_{i=1}^n i = \frac{n(n+1)}{2}
\end{equation}

\item 質量とエネルギーの等価性
\begin{equation}
E = mc ^2
\end{equation}

\item 波動方程式
\begin{equation}
\frac{\partial^2 z}{\partial t^2}=
c^2 (\frac{\partial^2 z}{\partial x^2}+\frac{\partial^2 z}{\partial y^2})-
\mu \frac{\partial z}{\partial t}
\end{equation}

\item 文章中に突然の $e^{i\theta}=\cos\theta+i\sin\theta$ オイラーの公式

\end{enumerate}
